\documentclass[10pt,a4paper]{amsart}
\usepackage{amsmath,amssymb,amsthm,graphicx}
\usepackage{verbatim}
\newenvironment{reference}{\comment}{\endcomment}
%\newenvironment{reference}{}{}
\newenvironment{slogan}{\comment}{\endcomment}
\newenvironment{history}{\comment}{\endcomment}
% For commutative diagrams we use Xy-pic
\usepackage[all]{xy}

% We use 2cell for 2-commutative diagrams.
\xyoption{2cell}
% Theorem environments.
%
\theoremstyle{plain}
\newtheorem{theorem}[subsection]{Theorem}
\newtheorem{proposition}[subsection]{Proposition}
\newtheorem{lemma}[subsection]{Lemma}

\theoremstyle{definition}
\newtheorem{definition}[subsection]{Definition}
\newtheorem{example}[subsection]{Example}
\newtheorem{exercise}[subsection]{Exercise}
\newtheorem{situation}[subsection]{Situation}

\theoremstyle{remark}
\newtheorem{remark}[subsection]{Remark}
\newtheorem{remarks}[subsection]{Remarks}

\numberwithin{equation}{subsection}

% Macros
%
\def\lim{\mathop{\mathrm{lim}}\nolimits}
\def\colim{\mathop{\mathrm{colim}}\nolimits}
\def\Spec{\mathop{\mathrm{Spec}}}
\def\Hom{\mathop{\mathrm{Hom}}\nolimits}
\def\Ext{\mathop{\mathrm{Ext}}\nolimits}
\def\SheafHom{\mathop{\mathcal{H}\!\mathit{om}}\nolimits}
\def\SheafExt{\mathop{\mathcal{E}\!\mathit{xt}}\nolimits}
\def\Sch{\mathit{Sch}}
\def\Mor{\mathop{\mathrm{Mor}}\nolimits}
\def\Ob{\mathop{\mathrm{Ob}}\nolimits}
\def\Sh{\mathop{\mathit{Sh}}\nolimits}
\def\NL{\mathop{N\!L}\nolimits}
\def\CH{\mathop{\mathrm{CH}}\nolimits}
\def\proetale{{pro\text{-}\acute{e}tale}}
\def\etale{{\acute{e}tale}}
\def\QCoh{\mathit{QCoh}}
\def\Ker{\mathop{\mathrm{Ker}}}
\def\Im{\mathop{\mathrm{Im}}}
\def\Coker{\mathop{\mathrm{Coker}}}
\def\Coim{\mathop{\mathrm{Coim}}}

% Boxtimes
%
\DeclareMathSymbol{\boxtimes}{\mathbin}{AMSa}{"02}

%
% Macros for moduli stacks/spaces
%
\def\QCohstack{\mathcal{QC}\!\mathit{oh}}
\def\Cohstack{\mathcal{C}\!\mathit{oh}}
\def\Spacesstack{\mathcal{S}\!\mathit{paces}}
\def\Quotfunctor{\mathrm{Quot}}
\def\Hilbfunctor{\mathrm{Hilb}}
\def\Curvesstack{\mathcal{C}\!\mathit{urves}}
\def\Polarizedstack{\mathcal{P}\!\mathit{olarized}}
\def\Complexesstack{\mathcal{C}\!\mathit{omplexes}}
% \Pic is the operator that assigns to X its picard group, usage \Pic(X)
% \Picardstack_{X/B} denotes the Picard stack of X over B
% \Picardfunctor_{X/B} denotes the Picard functor of X over B
\def\Pic{\mathop{\mathrm{Pic}}\nolimits}
\def\Picardstack{\mathcal{P}\!\mathit{ic}}
\def\Picardfunctor{\mathrm{Pic}}
\def\Deformationcategory{\mathcal{D}\!\mathit{ef}}
\usepackage{cleveref}
\crefname{originalequation}{original equation}{original equations}
\Crefname{originalequation}{Original equation}{Original equations}
\makeatletter
\newlabel{equation-towards-constructible-sets}{{equation-towards-constructible-sets}{source}{}{sourcepointer-1}{}}
\newlabel{equation-towards-constructible-sets@cref}{{[originalequation][0][]equation-towards-constructible-sets}{[1][0][]source}{}{}{}}
\makeatother
\crefname{originalreference}{original reference}{original references}
\Crefname{originalreference}{Original reference}{Original references}
\makeatletter
\newlabel{topology-lemma-spectral-inverse-limit-finite-sober-spaces}{{topology-lemma-spectral-inverse-limit-finite-sober-spaces}{source}{}{sourcepointer-2}{}}
\newlabel{topology-lemma-spectral-inverse-limit-finite-sober-spaces@cref}{{[originalreference][0][]topology-lemma-spectral-inverse-limit-finite-sober-spaces}{[1][0][]source}{}{}{}}
\makeatother
\crefname{originalreference}{original reference}{original references}
\Crefname{originalreference}{Original reference}{Original references}
\makeatletter
\newlabel{topology-lemma-directed-inverse-limit-spectral-spaces}{{topology-lemma-directed-inverse-limit-spectral-spaces}{source}{}{sourcepointer-3}{}}
\newlabel{topology-lemma-directed-inverse-limit-spectral-spaces@cref}{{[originalreference][0][]topology-lemma-directed-inverse-limit-spectral-spaces}{[1][0][]source}{}{}{}}
\makeatother
\crefname{originalreference}{original reference}{original references}
\Crefname{originalreference}{Original reference}{Original references}
\makeatletter
\newlabel{topology-lemma-descend-opens}{{topology-lemma-descend-opens}{source}{}{sourcepointer-4}{}}
\newlabel{topology-lemma-descend-opens@cref}{{[originalreference][0][]topology-lemma-descend-opens}{[1][0][]source}{}{}{}}
\makeatother
\crefname{originalreference}{original reference}{original references}
\Crefname{originalreference}{Original reference}{Original references}
\makeatletter
\newlabel{sheaves-lemma-compute-pullback-to-limit}{{sheaves-lemma-compute-pullback-to-limit}{source}{}{sourcepointer-5}{}}
\newlabel{sheaves-lemma-compute-pullback-to-limit@cref}{{[originalreference][0][]sheaves-lemma-compute-pullback-to-limit}{[1][0][]source}{}{}{}}
\makeatother
\title{A sheaf presented by finite data on a spectral space comes from a finite sober quotient}
\author{The Stacks Project Authors}
\date{stacks-project:modules.tex:local-snapshot-20261012; GNU FDL 1.2 or later}
\begin{document}
\maketitle
\noindent\textit{Editorial note.} All delivered mathematical text is the author's own, copied from the pinned Stacks Project chapter source sources/stacks-modules/pinned-original-source.tex; no statement or proof step was rewritten, reordered or removed.\par\medskip
\noindent\textit{Editorial note.} Preamble provenance: the Stacks Project keeps one single shared preamble for every chapter; this card takes it directly from the staged sources/stacks-modules/preamble.tex as three verbatim spans (lines 8-12, 14-18 and 154-217), so no derived file was written and nothing was added or redefined. The wrapper class and the base packages are supplied by the assembler.\par\medskip
\noindent\textit{Editorial note.} Bibliography provenance: neither the counted statement nor its proof contains a citation (this lane rescans the whole span before writing the card), so the card carries no bibliography and no entry was reconstructed.\par\medskip
\noindent\textit{Editorial note.} Selection provenance: this card was opened under the file-based control contract (CONTROL.md), with the mechanical pre-card gate run on its statement and proof interval before the card was written.\par\medskip
\noindent\textit{Allowed standard prerequisites (editorial).}\begin{itemize}
\item The presentation of a spectral space as a directed limit of finite sober spaces, with spectral transition maps and hence spectral projections. Invoked at the start of the counted proof (source lines 3289-3294) as cross-chapter pointers. Provides the finite quotients on which the presentation is to be descended.
\item The descent of opens and of maps into a colimit of sheaves along the projections of a directed limit. Invoked inside the counted proof (source lines 3295-3297 and 3316-3319) as cross-chapter pointers. Allows the finitely many opens and the finitely many values of the two maps to be brought down to some stage of the limit.
\end{itemize}\medskip
\ifdefined\hypertarget\hypertarget{sourcepointer-1}{}\fi
\noindent\textit{Source reference.} Original equation equation-towards-constructible-sets: sources/stacks-modules/pinned-original-source.tex, source lines 3281--3281. Navigation pointer only: the counted text refers to the Stacks item equation-towards-constructible-sets, which belongs to another chapter of the project and is not part of this card. The pointer carries the target's own label as its number because this chapter's staged snapshot does not contain the other chapter's tag table; its content is not reproduced here.\par\medskip
\ifdefined\hypertarget\hypertarget{sourcepointer-2}{}\fi
\noindent\textit{Source reference.} Original reference topology-lemma-spectral-inverse-limit-finite-sober-spaces: sources/stacks-modules/pinned-original-source.tex, source lines 3291--3291. Navigation pointer only: the counted text refers to the Stacks item topology-lemma-spectral-inverse-limit-finite-sober-spaces, which belongs to another chapter of the project and is not part of this card. The pointer carries the target's own label as its number because this chapter's staged snapshot does not contain the other chapter's tag table; its content is not reproduced here.\par\medskip
\ifdefined\hypertarget\hypertarget{sourcepointer-3}{}\fi
\noindent\textit{Source reference.} Original reference topology-lemma-directed-inverse-limit-spectral-spaces: sources/stacks-modules/pinned-original-source.tex, source lines 3293--3293. Navigation pointer only: the counted text refers to the Stacks item topology-lemma-directed-inverse-limit-spectral-spaces, which belongs to another chapter of the project and is not part of this card. The pointer carries the target's own label as its number because this chapter's staged snapshot does not contain the other chapter's tag table; its content is not reproduced here.\par\medskip
\ifdefined\hypertarget\hypertarget{sourcepointer-4}{}\fi
\noindent\textit{Source reference.} Original reference topology-lemma-descend-opens: sources/stacks-modules/pinned-original-source.tex, source lines 3297--3297. Navigation pointer only: the counted text refers to the Stacks item topology-lemma-descend-opens, which belongs to another chapter of the project and is not part of this card. The pointer carries the target's own label as its number because this chapter's staged snapshot does not contain the other chapter's tag table; its content is not reproduced here.\par\medskip
\ifdefined\hypertarget\hypertarget{sourcepointer-5}{}\fi
\noindent\textit{Source reference.} Original reference sheaves-lemma-compute-pullback-to-limit: sources/stacks-modules/pinned-original-source.tex, source lines 3317--3317. Navigation pointer only: the counted text refers to the Stacks item sheaves-lemma-compute-pullback-to-limit, which belongs to another chapter of the project and is not part of this card. The pointer carries the target's own label as its number because this chapter's staged snapshot does not contain the other chapter's tag table; its content is not reproduced here.\par\medskip
\begin{lemma}
\label{lemma-constructible-comes-from-finite}
Let $X$ be a spectral topological space. Let $\mathcal{B}$ be
the set of quasi-compact open subsets of $X$.
Let $\mathcal{F}$ be a sheaf of sets as in
Equation (\ref{equation-towards-constructible-sets}).
Then there exists a continuous spectral map $f : X \to Y$
to a finite sober topological space $Y$ and a sheaf
of sets $\mathcal{G}$ on $Y$ with finite stalks
such that $f^{-1}\mathcal{G} \cong \mathcal{F}$.
\end{lemma}

\begin{proof}
We can write $X = \lim X_i$ as a directed limit
of finite sober spaces, see Topology, Lemma
\ref{topology-lemma-spectral-inverse-limit-finite-sober-spaces}.
Of course the transition maps $X_{i'} \to X_i$ are spectral and hence
by Topology, Lemma \ref{topology-lemma-directed-inverse-limit-spectral-spaces}
the maps $p_i : X \to X_i$ are spectral.
For some $i$ we can find opens $U_{a, i}$ and $V_{b, i}$
of $X_i$ whose inverse images are $U_a$ and $V_b$, see
Topology, Lemma \ref{topology-lemma-descend-opens}.
The two maps
$$
\beta, \gamma :
\coprod\nolimits_{b \in B} j_{V_b!}\underline{S_b}
\longrightarrow
\coprod\nolimits_{a \in A} j_{U_a!}\underline{S_a}
$$
whose coequalizer is $\mathcal{F}$ correspond by adjunction to two families
$$
\beta_b, \gamma_b :
S_b
\longrightarrow
\Gamma(V_b, \coprod\nolimits_{a \in A} j_{U_a!}\underline{S_a}), \quad
b \in B
$$
of maps of sets. Observe that
$p_i^{-1}(j_{U_{a, i}!}\underline{S_a}) = j_{U_a!}\underline{S_a}$
and $(X_{i'} \to X_i)^{-1}(j_{U_{a, i}!}\underline{S_a}) =
j_{U_{a, i'}!}\underline{S_a}$. It follows from
Sheaves, Lemma \ref{sheaves-lemma-compute-pullback-to-limit}
(and using that $S_b$ and $B$ are finite sets) that
after increasing $i$ we find maps
$$
\beta_{b, i}, \gamma_{b, i} :
S_b
\longrightarrow
\Gamma(V_{b, i}, \coprod\nolimits_{a \in A} j_{U_{a, i}!}\underline{S_a})
, \quad b \in B
$$
which give rise to the maps $\beta_b$ and $\gamma_b$ after pulling
back by $p_i$. These maps correspond in turn to maps of sheaves
$$
\beta_i, \gamma_i :
\coprod\nolimits_{b \in B} j_{V_{b, i}!}\underline{S_b}
\longrightarrow
\coprod\nolimits_{a \in A} j_{U_{a, i}!}\underline{S_a}
$$
on $X_i$. Then we can take $Y = X_i$ and
$$
\mathcal{G} =
\text{Coequalizer}\left(
\xymatrix{
\coprod\nolimits_{b = 1, \ldots, m} j_{V_{b, i}!}\underline{S_b}
\ar@<1ex>[r] \ar@<-1ex>[r] &
\coprod\nolimits_{a = 1, \ldots, n} j_{U_{a, i}!}\underline{S_a}
}
\right)
$$
We omit some details.
\end{proof}

\end{document}
